% FortySecondsCV LaTeX template
% Copyright © 2019-2020 René Wirnata <rene.wirnata@pandascience.net>
% Licensed under the 3-Clause BSD License. See LICENSE file for details.
%
% Please visit https://github.com/PandaScience/FortySecondsCV for the most
% recent version! For bugs or feature requests, please open a new issue on
% github.
%
% Contributors
% ------------
% * ifokkema
% * Bertbk
% * Hespe
%
% Attributions
% ------------
% * fortysecondscv is based on the twentysecondcv class by Carmine Spagnuolo
%   (cspagnuolo@unisa.it), released under the MIT license and available under
%   https://github.com/spagnuolocarmine/TwentySecondsCurriculumVitae-LaTex
% * further attributions are indicated immediately before corresponding code


%-------------------------------------------------------------------------------
%                             ADDITIONAL PACKAGES
%-------------------------------------------------------------------------------
\documentclass[
    a4paper,
]{../fortysecondscv}


% improve word spacing and hyphenation
\usepackage{microtype}
\usepackage{ragged2e}

% Languages
\usepackage{fontspec}
\usepackage[english,greek]{babel}


% uncomment in case you don't want any hyphenation
% \usepackage[none]{hyphenat}

% take care of proper font encoding
\ifxetexorluatex
    \usepackage{fontspec}
    \defaultfontfeatures{Ligatures=TeX}
%    \newfontfamily\headingfont[Path = fonts/]{segoeuib.ttf} % local font
\else
    \usepackage[utf8]{inputenc}
    \usepackage[T1]{fontenc}
%    \usepackage[sfdefault]{noto} % use noto google font
\fi

% enable mathematical syntax for some symbols like \varnothing
\usepackage{amssymb}

% bubble diagram configuration
\usepackage{smartdiagram}
\smartdiagramset{
    % default font size is \large, so adjust to harmonize with sidebar layout
    bubble center node font = \footnotesize,
    bubble node font = \footnotesize,
    % default: 4cm/2.5cm; make minimum diameter relative to sidebar size
    bubble center node size = 0.4\sidebartextwidth,
    bubble node size = 0.25\sidebartextwidth,
    distance center/other bubbles = 1.5em,
    % set center bubble color
    bubble center node color = maincolor!70,
    % define the list of colors usable in the diagram
    set color list = {maincolor!10, maincolor!40,
    maincolor!20, maincolor!60, maincolor!35},
    % sets the opacity at which the bubbles are shown
    bubble fill opacity = 0.8,
}

% Bibliography
\usepackage[backend=biber,style=authoryear,sorting=ydnt]{biblatex}
\addbibresource{../bibs/my_publications.bib}

\usepackage{everypage}
\newcommand{\drawsidebar}{%
    \begin{tikzpicture}[remember picture,overlay]
        % fixed sidebar rectangle
        \fill[gray!10] (current page.north west) rectangle ([xshift=\sidebarwidth]current page.south west);


    \end{tikzpicture}%
}

%-------------------------------------------------------------------------------
%                            PERSONAL INFORMATION
%-------------------------------------------------------------------------------

%% mandatory information
% your name
\cvname{Ζαφειρόπουλος \\ Χάρης}

% job title/career
\cvjobtitle{
   Βιοπληροφορική,\\ Βιολογία Συστημάτων,\\ Μικροβιακή Οικολογία
}

% profile picture
\cvprofilepic{../pics/haris.jpeg}

% NOTE: ordering in sidebar will mimic the following order
% date of birth
\cvbirthday{Απρίλης 24, 1991}
% short address/location, use \newline if more than 1 line is required
\cvaddress{
   % Groefstraat 17, Leuven, Belgium
   Γαμβέττα 81, Θεσσαλονίκη, Ελλάδα
}
% phone number
\cvphone{+30 694 909 3089}
% personal website
\cvsite{https://hariszaf.github.io/}
% email address
\cvmail{haris.zafeiropoulos@atomicmail.io}
% pgp key
% \cvkey{4096R/FF00FF00}{0xAABBCCDDFF00FF00}

%-------------------------------------------------------------------------------
%                              SIDEBAR 1st PAGE
%-------------------------------------------------------------------------------
% add more profile sections to sidebar on first page
\addtofrontsidebar{
    % include gosquare national flags from https://github.com/gosquared/flags;
    % naming according to ISO 3166-1 alpha-2 country codes
    \graphicspath{{../pics/flags/}}

    \profilesection{Σχετικά με μένα}
    \aboutme{
      Αναπτύσσω υπολογιστικές και βιοπληροφορικές προσεγγίσεις για την καλύτερη κατανόηση των μικροβιακών κοινοτήτων, της δυναμικής τους και των αλληλεπιδράσεων μεταξύ των μικροοργανισμών. 
      Το κύριο επιστημονικό μου ενδιαφέρον εστιάζει στις μικροβιακές αλληλεπιδράσεις — ιδιαίτερα στις μεταβολικές — και στις εξελικτικές τους προεκτάσεις. Με ενδιαφέρει ιδιαίτερα ο τρόπος με τον οποίο οι μικροβιακές κοινότητες διαμορφώνονται από το περιβάλλον τους και, αντίστροφα, πώς η δομή και η λειτουργία τους επηρεάζουν το ίδιο το περιβάλλον. 
      Για τη διερεύνηση αυτών των ερωτημάτων, αξιοποιώ ένα ευρύ φάσμα μεθοδολογιών, συμπεριλαμβανομένων προσεγγίσεων «omics», ανάλυσης δικτύων και μεταβολικής μοντελοποίησης.
    }

    % social network accounts incl. proper hyperlinks
    \profilesection{Ερευνητικά Προφίλ}
        \begin{icontable}{2.5em}{1em}

            \social{\faGithub}
                {https://github.com/hariszaf}
                {Github Profile}

            \social{\faLinkedin}
                {https://www.linkedin.com/in/haris-zafeiropoulos-44555711a/}
                {LinkedIn Profile}

            \social{\faGraduationCap}
                {https://scholar.google.com/citations?user=RWsm3NYAAAAJ\&hl=en}
                {Google Scholar Profile}

            \social{\faResearchgate}
                {https://www.researchgate.net/profile/Haris-Zafeiropoulos}
                {ResearchGate Profile}

            % \social{\faSkype}
            %     {haris.zaf2}
            %     {haris.zaf2}

            % \social{\faTwitter}
            %     {https://twitter.com/haris_zaf}
            %     {@haris\_zaf}

            % \social{\faBlueSky}
            %     {https://orcid.org/0000-0002-1825-0097}
            %     {ORCID Profile}
        \end{icontable}


}
%-------------------------------------------------------------------------------
%                              SIDEBAR 2nd PAGE
%-------------------------------------------------------------------------------
\addtobacksidebar{

    % Languages
    % \profilesection{Γλώσσες}
    %     \pointskill{\flag{../pics/flags/GR.png}}{Greek}{5}
    %     \pointskill{\flag{../pics/flags/GB.png}}{English}{5}
    %     \pointskill{\flag{../pics/flags/FR.png}}{French}{3}


% Software
    \profilesection{Software}
    \begin{sidebarminipage}
        \chartlabel{\href{http://pema.hcmr.gr/}{PEMA}}
        \chartlabel{\href{http://prego.hcmr.gr/}{PREGO}}
        \chartlabel{\href{https://github.com/hariszaf/darn}{DARN}}
        \chartlabel{\href{https://github.com/GeomScale/dingo}{dingo}}
        \chartlabel{\href{https://hariszaf.github.io/microbetag}{microbetag}}
      \chartlabel{\href{https://github.com/MGXlab/DNNGIOR}{DNNGIOR}}
      \chartlabel{\href{https://mgrowthdb.gbiomed.kuleuven.be/}{μGrowthDB}}

    \end{sidebarminipage}

    % Research fields
    \profilesection{Research fields}
    \begin{figure}\centering
        \smartdiagram[bubble diagram]{
            \textcolor{white}{\textbf{microbial}} \\
            \textcolor{white}{\textbf{communities}},
            \textcolor{black!90}{data} \\
            \textcolor{black!90}{integration},
            \textcolor{black!90}{metagenomics},
            \textcolor{black!90}{metabolic} \\
            \textcolor{black!90}{modeling},
            \textcolor{black!90}{workflows} \\
            \textcolor{black!90}{\& HPC}
        }
    \end{figure}

    % Memberships
    \profilesection{Memberships}
        \begin{memberships}
            \membership[4em]{../pics/mkb.png}{\href{http://www.mikrobiokosmos.org/}{Mikrobiokosmos}}
            \membership[4em]{../pics/isme.png}{\href{https://www.isme-microbes.org/}{International Society of Microbial Ecology}}
            \membership[4em]{../pics/bsm.png}{\href{https://belsocmicrobio.be}{Belgian Society for Microbiology}}
        \end{memberships}
}

%-------------------------------------------------------------------------------
%                              SIDEBAR 3rd PAGE
%-------------------------------------------------------------------------------
\addtomoresidebar{

    \profilesection{Δεξιότητες \\ Προγραμματισμού}
        \pointskill{\faLinux}{Unix / Linux / AWK}{4}[5]
        \pointskill{\faPython}{Python 3.0}{4}[5]
        \pointskill{\flag{../pics/flags/Matlab_Logo.png}}{Matlab}{3}[5]
        \pointskill{\textsf{R}}{R}{3}[5]
        \pointskill{\flag{../pics/flags/cpp.png}}{C++}{2}[5]
        \pointskill{\faDev}{HTML / CSS / JS}{3}[5]
        \pointskill{\faDocker}{Docker \& Singularity}{4}[5]
        \pointskill{\faCloudDownload*}{REST API}{4}[5]


    % \profilesection{Math and Stats}
        % \pointskill{\faTable}{Γραμμική Άλγεβρα}{4}[5]
        % \pointskill{\faCube}{Mεταβολική Mοντελοποίηση βασισμένη σε Περιορισμούς}{4}[5]


    \profilesection{Εξειδικευμένες Ικανότητες}
        \barskill{\faVirus}{Μικροβιακή Οικολογία}{80}  % \faVirus
        \barskill{\faCubes}{Μοντελοποίηση Μεταβολισμού}{90}  % \faVirus
      \barskill{\faTable}{Γραμμική Άλγεβρα}{70}
        \barskill{\faLaptop}{Προγραμματισμός}{90} % \faLaptop
        \barskill{\faPenNib}{Συγγραφή Επιστημονικών Άρθρων}{90} % \faBookReader
}

%-------------------------------------------------------------------------------
%                              SIDEBAR EMPTY PAGES
%-------------------------------------------------------------------------------
\addtoemptysidebar{}



%-------------------------------------------------------------------------------
%                         TABLE ENTRIES RIGHT COLUMN
%-------------------------------------------------------------------------------
\begin{document}

\makefrontsidebar

\cvsection{Εργασιακή Εμπειρία}

\cvsubsection{Ερευνητικά έργα στο πλαίσιο επαγγελματικής απασχόλησης }

    \begin{cvtable}[1]

        \cvitem{2022 - 2026}{\href{https://3domics.eu}{3D'omics}}{Μεταδιδακτορικό}
        {
         \textbf{Στόχος του έργου} είναι η ανάπτυξη τεχνολογιών και πρωτοκόλλων για τη δημιουργία τρισδιάστατων «omics landscapes» και την ανακατασκευή οικοσυστημάτων αλληλεπίδρασης εντερικού μικροβιώματος και ξενιστή. 
         \textbf{Ο ρόλος μου} επικεντρώνεται στην ανάλυση μικροβιακών δικτύων, συνδυάζοντας χωρικά και μεταβολικά δεδομένα για την πρόβλεψη αλληλεπιδράσεων μεταξύ ειδών και της επίδρασης των μεταβολιτών στην κατανομή τους, μέσω ανάλυσης δικτύων συν-εμφάνισης και μεταβολικής μοντελοποίησης σε κλίμακα γονιδιώματος. 
        } 
      
      \\

        \cvitem{2022 - 2022}{
         \href{https://imbbc.hcmr.gr/project/eosc-life-gos/
      }{A workflow for marine Genomic Observatories data analysis}}{Επιστημονικός συνεργάτης} 
      {
            \textbf{Στόχος του έργου} ήταν η διευκόλυνση της ερμηνείας των μεγάλων όγκων δεδομένων που παράγονται από γονιδιωματικά παρατηρητήρια, μέσω της έγκαιρης παροχής των ταξινομικών χαρακτηρισμών κάθε δείγματος σε κατανοητή και προσβάσιμη μορφή. 
            \textbf{Ο ρόλος μου} ήταν η κύρια ανάπτυξη της υπολογιστικής ροής εργασίας που παραδόθηκε στο πλαίσιο του έργου, ενώ κατά το δεύτερο μισό του ανέλαβα καθήκοντα Επιστημονικά Υπεύθυνου (Principal Investigator).
        } 
      
      \\

        \cvitem{2018 - 2021}{\href{http://prego.hcmr.gr/}{PREGO: Process, environment, organism (PREGO)}}{Υποψήφιος Διδάκτορας}
            {
            \textbf{Στόχος του έργου} ήταν η μελέτη και κατανόηση της λειτουργίας των οικοσυστημάτων σε μικροβιακό επίπεδο, μέσω μιας προσέγγισης συστημικής βιολογίας. 
            \textbf{Ο ρόλος μου} στο έργο, το οποίο αποτέλεσε το κύριο ερευνητικό αντικείμενο της διδακτορικής μου διατριβής, ήταν η ανάπτυξη των επιμέρους λειτουργιών της πλατφόρμας που αφορούν την ολοκλήρωση και ενοποίηση ετερογενών βιολογικών δεδομένων.
        } 
      
      \\
        
        \cvitem{2019 -- present}{\href{https://elixir-greece.org}{ELIXIR–GR}}{Επιστημονικός Συνεργάτης}
            {
         \textbf{Στόχος του έργου} ήταν η συμβολή στην ανάπτυξη μιας βιώσιμης ευρωπαϊκής υποδομής για βιολογικές πληροφορίες, στο πλαίσιο του ELIXIR-GR, του Ελληνικού Εθνικού Κόμβου της ευρωπαϊκής ερευνητικής υποδομής ELIXIR. 
         \textbf{Ο ρόλος μου} στο έργο ήταν η ανάπτυξη containerized εκδόσεων διαφόρων εργαλείων λογισμικού, με στόχο τη διευκόλυνση της εγκατάστασης, αναπαραγωγιμότητας και αξιοποίησής τους.   
        }

      \\
        
        \cvitem{2018 - 2020}{\href{https://reconnect.hcmr.gr/}{RECONNECT}}{Επιστημονικός Συνεργάτης}
            {
         \textbf{Στόχος του έργου} (Interreg V-B «Balkan-Mediterranean 2014–2020»)η ανάπτυξη στρατηγικών για τη βιώσιμη διαχείριση Θαλάσσιων Προστατευόμενων Περιοχών (MPAs) και περιοχών Natura 2000. 
         \textbf{Ο ρόλος μου} επικεντρώθηκε στην ανάλυση μικροβιακών δεδομένων αλληλούχισης  (amplicon analysis).
        }

    \end{cvtable}



\cvsubsection{Συμμετοχές σε άλλα έργα}

    \begin{cvtable}[1]

        \cvitem{2021 - 2023}{\href{https://www.embrc.eu}{EMBRC}}{Επιστημονικός Συνεργάτης}
            {
         \textbf{Στόχος του έργου} ήταν η υποστήριξη της βασικής και εφαρμοσμένης έρευνας στη θαλάσσια βιολογία και οικολογία, καθώς και η προώθηση μιας βιώσιμης «γαλάζιας οικονομίας», μέσω του European Marine Biological Resource Centre (EMBRC). 
         \textbf{Ο ρόλος μου} στο έργο ήταν να συνεισφέρω στις δραστηριότητες του 
         \href{https://www.embrc.eu/embrc-network/greece}{Ελληνικού Κόμβου του EMBRC (CMBR)}, με έμφαση στην ανάπτυξη υπολογιστικών ροών εργασίας (workflows) για την ανάλυση και επεξεργασία των δεδομένων που παράγονται από το δίκτυο.
        }

     \\

     \cvitem{2020 - ongoing}
         {\href{https://www.lifewatch.eu}{LifeWatch-ERIC}}
         {Επιστημονικός Συνεργάτης}
         {
            \textbf{Στόχος του έργου:} Η μελέτη και μακροχρόνια παρακολούθηση της βιοποικιλότητας, με έμφαση στα εισβολικά θαλάσσια είδη, στο πλαίσιο του LifeWatch ERIC.
            \textbf{Ο ρόλος μου:} Συντήρηση και ανάπτυξη του workflow του PEMA σύμφωνα με τις ανάγκες της υποδομής, στο πλαίσιο της~\href{https://www.lifewatch.eu/internal-joint-initiative/validation-cases/long-term-monitoring-of-invasive-marine-species/}{Internal Joint Initiative για την παρακολούθηση εισβλητικών ειδών}.
         }

    \end{cvtable}


%  -----------------------------------------------------------------------
% \newpage
\makemoresidebar

%  -----------------------------------------------------------------------
\AddEverypageHook{\drawsidebar}

% ----------------------------------------------------------------------------
%                            EDUCATION
% ----------------------------------------------------------------------------

\cvsection{Εκπαίδευση}

    \cvsubsection{Μεταπτυχιακές Σπουδές}

    \begin{cvtable}[1]
        \cvitem{2018 -- 2022}{ Διδακτορικό στη Βιοπληροφορική }
        {\href{https://en.uoc.gr/}{Πανεπιστήμιο Κρήτης}, \href{https://www.biology.uoc.gr/en}{Τμήμα Βιολογίας}}{}
        \cvitem{}{Διατριβή}{\href{https://imbbc.hcmr.gr/}{ΙΘΑΒΒΥΚ - ΕΛΚΕΘΕ}}
            {\href{https://www.openarchives.gr/aggregator-openarchives/edm/elocus/000018-dlib_5_7_d_metadata-dlib-1661248925-407543-18532.tkl?language=en}{Μικροβιακές κοινότητες υπό το πρίσμα αλληλουχίας υψηλής απόδοσης, ολοκλήρωσης δεδομένων \& ανάλυση μεταβολικών δικτύων}}

      \\
 
      \cvitem{2016 -- 2018}{Μεταπτυχιακό στη Βιοπληροφορική}
        {\href{https://en.uoc.gr/}{Πανεπιστήμιο Κρήτης}, \href{http://www.english.med.uoc.gr/}{Ιατρική Σχολή}}
            {Βαθμός: 9.1/10.0}
        \cvitem{}{Πτυχιακή εργασία}{\href{https://imbbc.hcmr.gr/}{ΙΘΑΒΒΥΚ - ΕΛΚΕΘΕ}}
            {\href{https://www.openarchives.gr/aggregator-openarchives/edm/elocus/000018-dlib_5_2_f_metadata-dlib-1545038085-364284-26948.tkl}{Μετακωδικοποίηση eDNA για την αξιολόγηση της βιοποικιλότητας: σχεδιασμός αλγορίθμων και υλοποίηση γραμμών εργασιών βιοπληροφορικής}}

    \end{cvtable}


    \cvsubsection{Προπτυχιακές Σπουδές}
        \begin{cvtable}[1]
            
        \cvitem{2011 -- 2016}{Πτυχίο στη Βιολογία}{\href{https://en.uoa.gr/}{Εθνικό \& Καποδιστριακό Πανεπιστήμιο Αθηνών} }
            {Βαθμός: 6.2 / 10.0}
        \cvitem{}{Πτυχιακή εργασία}{\href{https://en.uoa.gr/schools_and_departments/school_of_science/}{Σχολή Θετικών Επιστημών}, \href{http://en.biol.uoa.gr/}{Τμήμα Βιολογίας}}
            {Μορφολογία, μορφομετρία και ανατομία ειδών του γένους \textit{Pseudamnicola} στην Ελλάδα}
            
    \end{cvtable}


% % ----------------------------------------------------------------------------
% %                            PUBLICATIONS
% % ----------------------------------------------------------------------------
\cvsection{Δημοσιεύσεις}

   \begin{cvtable}
       \cvpubitem{Sampling from the Solution Space and Metabolic Environments of Genome-Scale Metabolic Models}{\textbf{Zafeiropoulos, Haris} and Rios Garza, Daniel}{arXiv, DOI: \href{https://doi.org/10.48550/arXiv.2603.29546}{10.48550/arXiv.2603.29546}}{2026}
    \cvpubitem{microbetag: simplifying microbial network interpretation through annotation, enrichment tests, and metabolic complementarity analysis}{\textbf{Zafeiropoulos, Haris}, Delopoulos, Ermis Ioannis Michail, Erega, Andi, Schneider, Aline, Geirnaert, Annelies, Morris, John and Faust, Karoline}{Genome Biology, DOI: \href{https://doi.org/10.1186/s13059-025-03769-2}{10.1186/s13059-025-03769-2}}{2025}
    \cvpubitem{Establishing the ELIXIR Microbiome Community}{Finn, R. D., Balech, B., Burgin, J., Chua, P., Corre, E., Cox, C. J., Donati, C., dos Santos, V. M., Fosso, B., Hancock, J., Heil, K. F., Ishaque, N., Kale, V., Kunath, B. J., Médigue, C., Nogueira, T., Pafilis, Evangelos, Pesole, G., Richardson, L., Santamaria, M., Strepis, N., Van Den Bossche, T., Vizcaíno, J. A., \textbf{Zafeiropoulos, Haris}, Willassen, N. P., Pelletier, E. and Batut, B.}{F1000Research, DOI: \href{https://doi.org/10.12688/f1000research.144515.2}{10.12688/f1000research.144515.2}}{2025}
    \cvpubitem{\ensuremath{\mu}GrowthDB: querying, visualizing, and sharing microbial growth curve data}{Radev, Andrey, Casado Gómez-Pallete, Julia, Monsalve Duarte, Sofia, Faust, Karoline and \textbf{Zafeiropoulos, Haris}}{bioRxiv, DOI: \href{https://doi.org/10.1101/2025.10.13.682118}{10.1101/2025.10.13.682118}}{2025}
    \cvpubitem{A long-term ecological research data set from the Marine Genetic Monitoring Program ARMS-MBON 2018--2020}{Daraghmeh, Nauras, Exter, Katrina, Pagnier, Justine, Balazy, Piotr, Cancio, Ibon, Chatzigeorgiou, Giorgos, Chatzinikolaou, Eva, Chelchowski, Maciej, Chrismas, Nathan Alexis Mitchell, Comtet, Thierry and others}{Molecular Ecology Resources, DOI: \href{https://doi.org/10.1111/1755-0998.14073}{10.1111/1755-0998.14073}}{2025}
    \cvpubitem{Predicting microbial interactions with approaches based on flux balance analysis: an evaluation}{Joseph, Clémence, \textbf{Zafeiropoulos, Haris}, Bernaerts, Kristel and Faust, Karoline}{BMC Bioinformatics, DOI: \href{https://doi.org/10.1186/s12859-024-05651-7}{10.1186/s12859-024-05651-7}}{2024}
    \cvpubitem{Improving genome-scale metabolic models of incomplete genomes with deep learning}{Boer, Meine D., Melkonian, Chrats, \textbf{Zafeiropoulos, Haris}, Haas, Andreas F., Rios Garza, Daniel and Dutilh, Bas E.}{iScience, DOI: \href{https://doi.org/10.1016/j.isci.2024.111349}{10.1016/j.isci.2024.111349}}{2024}
    \cvpubitem{dingo: a Python package for metabolic flux sampling}{Chalkis, Apostolos, Fisikopoulos, Vissarion, Tsigaridas, Elias and \textbf{Zafeiropoulos, Haris}}{Bioinformatics Advances, DOI: \href{https://doi.org/10.1093/bioadv/vbae037}{10.1093/bioadv/vbae037}}{2024}
    \cvpubitem{Unveiling emerging opportunistic fish pathogens in aquaculture: a comprehensive seasonal study of microbial composition in Mediterranean fish hatcheries}{Skliros, Dimitrios, Kostakou, Maria, Kokkari, Constantina, Tsertou, Maria Ioanna, Pavloudi, Christina, \textbf{Zafeiropoulos, Haris}, Katharios, Pantelis and Flemetakis, Emmanouil}{Microorganisms, DOI: \href{https://doi.org/10.3390/microorganisms12112281}{10.3390/microorganisms12112281}}{2024}
    \cvpubitem{metaGOflow: a workflow for the analysis of marine Genomic Observatories shotgun metagenomics data}{\textbf{Zafeiropoulos, Haris}, Beracochea, Martin, Ninidakis, Stelios, Exter, Katrina, Potirakis, Antonis, De Moro, Gianluca, Richardson, Lorna, Corre, Erwan, Machado, João, Pafilis, Evangelos and others}{GigaScience, DOI: \href{https://doi.org/10.1093/gigascience/giad078}{10.1093/gigascience/giad078}}{2023}
    \cvpubitem{Metabolic models of human gut microbiota: advances and challenges}{Rios Garza, Daniel, Gonze, Didier, \textbf{Zafeiropoulos, Haris}, Liu, Bin and Faust, Karoline}{Cell Systems, DOI: \href{https://doi.org/10.1016/j.cels.2022.11.002}{10.1016/j.cels.2022.11.002}}{2023}
    \cvpubitem{PREGO: a literature and data-mining resource to associate microorganisms, biological processes, and environment types}{\textbf{Zafeiropoulos, Haris}, Paragkamian, Savvas, Ninidakis, Stelios, Pavlopoulos, Georgios A., Jensen, Lars Juhl and Pafilis, Evangelos}{Microorganisms, DOI: \href{https://doi.org/10.3390/microorganisms10020293}{10.3390/microorganisms10020293}}{2022}
    \cvpubitem{Deciphering the community structure and the functional potential of a hypersaline marsh microbial mat community}{Pavloudi, Christina and \textbf{Zafeiropoulos, Haris}}{FEMS Microbiology Ecology, DOI: \href{https://doi.org/10.1093/femsec/fiac141}{10.1093/femsec/fiac141}}{2022}
    \cvpubitem{Automating the curation process of historical literature on marine biodiversity using text mining: the DECO workflow}{Paragkamian, Savvas, Sarafidou, Georgia, Mavraki, Dimitra, Pavloudi, Christina, Beja, Joana, Eliezer, Menashé, Lipizer, Marina, Boicenco, Laura, Vandepitte, Leen, Perez-Perez, Ruben and others}{Frontiers in Marine Science, DOI: \href{https://doi.org/10.3389/fmars.2022.940844}{10.3389/fmars.2022.940844}}{2022}
    \cvpubitem{0s and 1s in marine molecular research: a regional HPC perspective}{\textbf{Zafeiropoulos, Haris}, Gioti, Anastasia, Ninidakis, Stelios, Potirakis, Antonis, Paragkamian, Savvas, Angelova, Nelina, Antoniou, Aglaia, Danis, Theodoros, Kaitetzidou, Eliza, Kasapidis, Panagiotis and others}{GigaScience, DOI: \href{https://doi.org/10.1093/gigascience/giab053}{10.1093/gigascience/giab053}}{2021}
    \cvpubitem{The Dark mAtteR iNvestigator (DARN) tool: getting to know the known unknowns in COI amplicon data}{\textbf{Zafeiropoulos, Haris}, Gargan, Laura, Hintikka, Sanni, Pavloudi, Christina and Carlsson, Jens}{Metabarcoding and Metagenomics, DOI: \href{https://doi.org/10.3897/mbmg.5.69657}{10.3897/mbmg.5.69657}}{2021}
    \cvpubitem{Geometric Algorithms for Sampling the Flux Space of Metabolic Networks}{Chalkis, Apostolos, Fisikopoulos, Vissarion, Tsigaridas, Elias and \textbf{Zafeiropoulos, Haris}}{37th International Symposium on Computational Geometry (SoCG 2021), DOI: \href{https://doi.org/10.4230/LIPIcs.SoCG.2021.21}{10.4230/LIPIcs.SoCG.2021.21}}{2021}
    \cvpubitem{The Santorini Volcanic Complex as a valuable source of enzymes for bioenergy}{Polymenakou, Paraskevi N., Nomikou, Paraskevi, \textbf{Zafeiropoulos, Haris}, Mandalakis, Manolis, Anastasiou, Thekla I., Kilias, Stephanos, Kyrpides, Nikos C., Kotoulas, Georgios and Magoulas, Antonios}{Energies, DOI: \href{https://doi.org/10.3390/en14051414}{10.3390/en14051414}}{2021}
    \cvpubitem{PEMA: a flexible pipeline for environmental DNA metabarcoding analysis of the 16S/18S ribosomal RNA, ITS, and COI marker genes}{\textbf{Zafeiropoulos, Haris}, Viet, Ha Quoc, Vasileiadou, Katerina, Potirakis, Antonis, Arvanitidis, Christos, Topalis, Pantelis, Pavloudi, Christina and Pafilis, Evangelos}{GigaScience, DOI: \href{https://doi.org/10.1093/gigascience/giaa022}{10.1093/gigascience/giaa022}}{2020}

   \end{cvtable}

% ------------------------------------------------------------------------
%                    CONFERENCES 
% ------------------------------------------------------------------------
\cvsection{Συνέδρια}

   \begin{cvtable}[1]
          \cvitem{2025}{Poster presentation: {\textbackslash}texttt\{microbetag\}: simplifying microbial network interpretation through annotation, enrichment and metabolic complementarity analysis}{\href{https://icsb2025.com/}{ICSB25}}{24th Annual International Conference on Systems Biology (ICSB 2025) @ Dublin, Ireland}{}\\

    \cvitem{2024}{Poster presentation: {\textbackslash}texttt\{microbetag\}: simplifying microbial network interpretation through annotation, enrichment and metabolic complementarity analysis}{\href{https://isme-microbes.org/isme-symposia/past-isme-symposia/}{ISME19}}{19th International Symposium on Microbial Ecology @ Cape Town, South Africa}{}\\

    \cvitem{2024}{Poster presentation: {\textbackslash}texttt\{microbetag\}: simplifying microbial network interpretation through annotation, enrichment and metabolic complementarity analysis}{\href{https://belsocmicrobio.be/events/annual-symposium-2024/}{BSM2024}}{Belgian Society for Microbiology: Annual Symposium @ Brussels, Belgium}{}\\

    \cvitem{2022}{Poster presentation: {\textbackslash}texttt\{microbetag\}: a microbial co-occurrence network annotator}{\href{https://appliedhologenomicsconference.eu}{AHC2022}}{1st Applied HoloGenomics conference @ Bilbao, Spain}{}\\

    \cvitem{2022}{Poster presentation: Reverse ecology and systems biology approaches to identify key processes ensuring life in extreme environments}{\href{https://isme18.isme-microbes.org}{ISME18}}{18th International Symposium on Microbial Ecology @ Lausanne, Switzerland}{}\\

    \cvitem{2021}{Flash talk: {\textbackslash}texttt\{dingo\}: A python library for metabolic networks sampling {\textbackslash}\& analysis}{\href{https://www.open-bio.org/events/bosc-2021/}{BOSC2021}}{Bioinformatics Open Source Conference (BOSC) - online}{}\\

    \cvitem{2021}{Flash talk: PEMA v2: addressing metabarcoding bioinformatics analysis challenges}{\href{https://symposium.inrae.fr/dnaqua-conference-evian2021/}{DNAQUA2021}}{1st DNAQUA International Conference - online}{}\\

    \cvitem{2020}{Flash talk: Mining literature and -omics (meta)data to associate microorganisms, biological processes and environment types}{\href{https://fems2020belgrade.com/}{FEMS2020}}{Federation of European Microbiological Societies (FEMS) - online}{}\\

    \cvitem{2020}{Oral presentation: Geometric and statistical methods in systems biology: the case of metabolic networks}{\href{https://pydata.org/global2020/}{PyData2020}}{PyData Global - online}{}\\

    \cvitem{2019}{Participation}{\href{https://network2019.sciencesconf.org}{Network2019}}{network2019: 4th Symposium on Ecological Networks @ Paris, France}{}\\

    \cvitem{2019}{ePoster presentation: P.E.M.A.: a Pipeline for Environmental DNA Metabarcoding Analysis}{\href{http://dnabarcodes2019.org/}{iBOL2019}}{8th International Barcode of Life Conference @ Trondheim, Norway}{}\\

    \cvitem{2018}{Poster presentation: P.E.M.A.: a Pipeline for Environmental DNA Metabarcoding Analysis}{\href{https://www.iscb.org/cms_addon/events/details.php?uid=2508}{ECCB18}}{European Conference on Computational Biology (ECCB) 2018 @ Athens, Greece}{}\\


   \end{cvtable}


% ------------------------------------------------------------------------
%                    SUMMER SCHOOLS
% ------------------------------------------------------------------------
\cvsection{Εργαστήρια \& Summer schools}

   \begin{cvtable}[1]
      

    \cvitem{2025}{School of artificial intelligence applied to microbiomes @ AgroParisTech, France}
        {\href{https://ai-microbiome-school.onrender.com/}{SAIM}}
        {
            - Sampling the solution space of metabolic networks (theory \& hands-on)
        }{}

    \\[0.2ex]

    \cvitem{2025}{Data integration with Microbial Networks and Community Models @ KULeuven, Belgium}
        {\href{http://msysbiology.com/microbialdataintegration/}{DIMNCM}}
        {
            - \textit{Organizing committee}
            -  microbetag: metabolic secrets behind microbial co-occurrence (theory \& hands-on)
        }{}

    \\[0.2ex]


    \cvitem{2025}{Metabolic models applied to microbiomes @ Antony, France}
        {\href{https://metabolicmodelingantony2025.onrender.com/}{MMAM}}
        {
            - \textit{Organizing committee}
            - Metabolic model analysis methods (theory \& hands-on)
        }{}

    \\[0.2ex]
    
    \cvitem{2023}{Summer School on Microbial Community Design and Control @ Leuven, Belgium}
        {\href{http://msysbiology.com/microbedesign/}{DCMC}}
        {Assisting on "metabolic network reconstruction" topic}{}

    \\[0.2ex]

    \cvitem{2023}{Economic Principles in Cell Physiology @ Paris, France}
        {\href{https://appliedhologenomicsconference.eu}{EPCP2023}}
        {Mathematical modeling of cellular systems and “resource allocation thinking”. 
        \textit{Poster presentation: \\[0.2ex] "sampl'em all"}: exploring potential flux sampling applications}{}
    \\[0.2ex]
 
    \cvitem{2022}{Microbial communities: current approaches and open challenges @ Cambridge, UK}
        {\href{https://www.newton.ac.uk/event/umcw06/}{UMCW06}}
        {Interdisciplinary exchanges in microbial communities' research and the emerging challenge areas.}{}
    \\[0.2ex]

    \cvitem{2022}{ELIXIR Fluxomics Training School 2021}
    {\href{https://tess.elixir-europe.org/events/elixir-fluxomics-training-school-2021}{ELIXFLUX}}
    {An introduction to the field of fluxomics and the experimental and computational methods used to estimate and predict metabolic fluxes}{}

    \\[0.2ex]

    \cvitem{2020}{ Metagenomics, Metatranscript- omics and multi 'omics for microbial community studies}
    {\href{https://docs.google.com/spreadsheets/d/1YwJWejTxh7FLfooKCEJ18QCqvkYKAe5jFtwqQ7qU7bI/edit?gid=674784997\#gid=674784997}{material}}
    {A Physalia Course from Prof. Dr. Curtis Huttenhower}{}


   \end{cvtable}

% ------------------------------------------------------------------------
%                    SUPERVISION
% ------------------------------------------------------------------------
\cvsection{Επίβλεψη Φοιτητών} 

   \begin{cvtable}[1]
   \cvitem{2025--ongoing}{Sourav Pattanaik, MSc student}
{{}}
{"Populating $\mu$GgrowrhDB"}
{}
\\

\cvitem{2024--ongoing}{Sotiris Touliopoulos, ERASMUS student}
{\href{https://github.com/SotirisTouliopoulos/dingo-stats}{GitHub repo}}
{"Pre- and post-sampling features to leverage flux sampling at both the strain and the community level"}
{}
\\

\cvitem{2024--ongoing}{Xi Peng, Visiting PhD candidate}
{{}}
{"Enhancing microbetag's functional annotations with quorum sensing prediction capabilities"}
{}
\\

\cvitem{2024--ongoing}{Anna Voukouna, MSc student}
{{}}
{"Investigation of microbial metabolic interactions potential across the Genome Taxonomy Database representative genomes"}
{}
\\

\cvitem{2024--2025}{Andrey Radev, MSc student}
{{}}
{"Visualisation and analysis of microbial growth data to facilitate interactions investigation and interactions' strength assessment"}
{}
\\

\cvitem{2024}{Sotiris Touliopoulos, GSoC student}
{\href{https://sotiristouliopoulos.github.io/dingo/}{blogpost}}
{"Pre- and post-sampling features to leverage flux sampling at both the strain and the community level"}
{}
\\

\cvitem{2023--2024}{Sofia Monsalve Duarte, MSc student}
{{}}
{"Creation and implementation of a database interface for the collection and analysis of microbial growth data"}
{}
\\

\cvitem{2023--2024}{Ermis Ioannis Michail Delopoulos, GSoC student}
{\href{https://ermismd.github.io/MGG/}{post}}
{"Development of a Cytoscape App for microbe-microbe association networks"}
{}
\\

\cvitem{2022--2024}{Ermis Ioannis Michail Delopoulos, MSc student}
{{}}
{"Development of a graphical user interface for microbetag"}
{}
\\

\cvitem{2022--2023}{Julia Casado Gómez-Pallete, MSc student}
{\href{https://kuleuven.limo.libis.be/discovery/fulldisplay?docid=alma9993527122501488&context=L&vid=32KUL_KUL:KULeuven&search_scope=All_Content&tab=all_content_tab&lang=en}{thesis}}
{"Development of a database for human gut bacterial growth curves"}
{}
\\


   \end{cvtable}

% ------------------------------------------------------------------------
%                    TEACHING 
% ------------------------------------------------------------------------
\cvsection{Εκπαιδευτική Εμπειρία}

   \begin{cvtable}[1]
          \cvitem{2025}{Exploring the microbiome with high-throughput sequencing technologies: a bioinformatics perspective}{\href{https://docs.google.com/presentation/d/14HP6WUIppr5tlEridW6_CqpBHKRtNS1JJJCaKguLC6A/edit?usp=sharing}{slides}}{Lecture on MSc "Applied Bioinformatics \& Data Analysis" in \href{https://bioinfo.mbg.duth.gr}{Democritus University of Thrace}, Greece}{}\\
    \cvitem{2025}{Introduction in metabolic modeling}{\href{https://github.com/hariszaf/metabolic_toy_model/tree/duth}{GitHub repo}}{Theoretical and practical session on MSc "Applied Bioinformatics \& Data Analysis" in \href{https://bioinfo.mbg.duth.gr}{Democritus University of Thrace}, Greece}{}\\
    \cvitem{2024}{Exploring the microbiome with high-throughput sequencing technologies: a bioinformatics perspective (part II)}{\href{https://docs.google.com/presentation/d/1jNNGnUWEOk40eIGnJemGaThLcbmtyydaWyJZtGpLj7w/edit?usp=sharing}{slides}}{Lecture on MSc "Applied Bioinformatics \& Data Analysis" in \href{https://bioinfo.mbg.duth.gr}{Democritus University of Thrace}, Greece}{}\\
    \cvitem{2024}{Exploring the microbiome with high-throughput sequencing technologies: a bioinformatics perspective}{\href{https://docs.google.com/presentation/d/1ur_r9DRvZInixYkgDDYZNmVHH9mr15BKlIoVoCiyzhc/edit?usp=sharing}{slides}}{Lecture on MSc "Applied Bioinformatics \& Data Analysis" in \href{https://bioinfo.mbg.duth.gr}{Democritus University of Thrace}, Greece}{}\\
    \cvitem{2023}{eDNA metabarcoding, pipeline development \& high performance computing}{\href{https://docs.google.com/presentation/d/1FSObXaOt0HFq30UBypO4ui13oy1cbflxLRQJPQwfjC4/edit?usp=sharing}{slides}}{Lecture on MSc "Applied Bioinformatics \& Data Analysis" in \href{https://bioinfo.mbg.duth.gr}{Democritus University of Thrace}, Greece}{}\\
    \cvitem{2023}{eDNA metabarcoding, pipeline development \& high performance computing}{\href{https://docs.google.com/presentation/d/1FSObXaOt0HFq30UBypO4ui13oy1cbflxLRQJPQwfjC4/edit?usp=sharing}{slides}}{Lecture on MSc "Applied Bioinformatics \& Data Analysis" in \href{https://bioinfo.mbg.duth.gr}{Democritus University of Thrace}, Greece}{}\\
    \cvitem{2022}{Genome-scale model reconstruction}{\href{https://colab.research.google.com/drive/1uL-oiSNAQoAL1qWyS4z_y94kWhHyT4hA?usp=sharing}{GColab notebook}}{Lectures on "Master en bioinformatique et modélisation" of the \href{https://www.ulb.be/fr/programme/ma-binf}{Université Libre de Bruxelles}}{}\\
    \cvitem{2021}{Microbiome Data Analyses Workshop Online}{\href{https://meetinghand.com/e/mdawo}{MDAWO}}{\textbf{Tutorial:} PEMA: a flexible Pipeline for Environmental DNA Metabarcoding Analysis of the 16S/18S rRNA, ITS and COI marker genes.}{}\\
    \cvitem{2020}{Environmental DNA \& DNA metabarcoding}{\href{https://docs.google.com/presentation/d/19TVFO3Rw-1owvrqVdzmUzQUrYHZ1Vaw4ctVjBso-n-E/edit?usp=sharing}{slides}}{Lecture on MSc "Environmental Biology" in \href{http://envbio.biology.uoc.gr}{University of Crete}}{}\\
    \cvitem{2019}{8th Conference of Hellenic Society of Mikrobiokosmos}{\href{https://mikrobiokosmos.org/wp-content/uploads/2021/10/8th_mbk_abstract_book_2019.pdf}{MKB8}}{\textbf{Tutorial:} PEMA: a flexible Pipeline for Environmental DNA Metabarcoding Analysis of the 16S/18S rRNA, ITS and COI marker genes.}{}\\

   \end{cvtable}



% ------------------------------------------------------------------------
%                           AWARDS & REFS
% ------------------------------------------------------------------------
\cvsection{Βραβεία \& Υποτροφίες}

   \begin{cvtable}[1]
      
    \cvitem{2021}{9th conference of the Hellenic Society of Mikrobiokosmos}
    {\href{https://mikrobiokosmos.org/wp-content/uploads/2022/03/9th_mbk_abstract_book_2021.pdf}{MBK9}}
    {
        \textbf{Oral presentation award:}{
        Sampling the flux space of microbial metabolic networks for (un)biased assessment of the possible functional states:
        the example of Sars-Cov-2 on the human alveolar macrophage metabolic network.
    }}{}

    \\[0.2ex]

    \cvitem{2021}{Google Summer of Code}
    {\href{https://summerofcode.withgoogle.com/}{GSoC}}
    {Project title: \textit{\href{https://hariszaf.github.io/gsoc2021/}
    {From DNA sequences to metabolic interactions: building a pipeline to extract key metabolic processes}}}{}

    \\[0.2ex]

    \cvitem{2021}{European Molecular Biology Organization Short-Term Fellowship}
    {\href{https://www.embo.org/}{EMBO}}
    {Project title: \textit{"Exploiting data integration, text-mining and computational geometry to enhance microbial interactions inference from  co-occurrence networks"}}{}

    \\[0.2ex]

    \cvitem{2020}{Federation of European Microbiological Societies \\ Meeting Attendance Grant}
    {\href{https://fems-microbiology.org/}{FEMS}}
    {for joining the \textit{"Metagenomics, Metatranscript- omics and \\ 
    multi ’omics for microbial community studies"} Physalia course}{}

    \\[0.2ex]
    
    \cvitem{2019}{Short Term Scientific Mission (STSM) \\ @ DNAqua-net COST action}
    {\href{http://dnaqua.net/}{DNAqua-net}}
    {Project title \textit{A comparison of bioinformatic pipelines and sampling techniques to enable benchmarking of DNA metabarcoding}. \href{http://dnaqua.net/wp-content/uploads/2019/08/Zafeiropoulos.pdf}{Report.}}{}


    \\[0.2ex]
    
    \cvitem{2018}{Best Poster Award @ Hellenic Bioinformatics conference}
    {\href{https://hscbio.wordpress.com/}{HBIO}}
    {for \textit{PEMA: a Pipeline for Environmental DNA Metabarcoding Analysis}}{}
   \end{cvtable}


% ------------------------------------------------------------------------
%                           REFERENCES
% ------------------------------------------------------------------------
\cvsection{Συστάσεις}

   \begin{cvtable}[1]
      
    \cvitemshort{\href{http://msysbiology.com}{Karoline Faust}}
        {Department of Microbiology, Immunology and Transplantation, Rega Institute for Medical Research, Laboratory of Molecular Bacteriology, KU Leuven,
        Leuven, 3000, Leuven, Belgium
        \\ 
        \faPhone: +3216322698 | 
        \faAt: \href{mailto:karoline.faust@kuleuven.be}{karoline.faust@kuleuven.be}}
      
    \\ 


    % \cvitemshort{\href{http://lab42open.hcmr.gr/people/evangelospafilis/}{Ευάγγελος Παφίλης}}
    %     {Ινστιτούτο Θαλάσσιας Βιολογίας, Βιοτεχνολογίας \& Ιχθυοκαλλιεργιών, Ελληνικό Κέντρο Θαλάσσιας Βιολογίας, (ΙΘΑΒΒΥΚ - ΕΛΚΕΘΕ),
    %     Γούρνες, Πεδιάδος, Ηράκλειο, Κρήτη 71003, Ελλάδα \\ \faPhone: +30 2810 337740 | 
    %     \faAt: \href{mailto:pafilis@hcmr.gr}{pafilis@hcmr.gr}}
        
    % \\


    \cvitemshort{\href{https://m3i-lab.github.io/people/johannes.html}{Johanens Zimmermann}}
        {Cluster of Excellence Balance of the Microverse,
        Friedrich Schiller University Jena 
        Rosalind-Franklin-Strasse 1, 07745 Jena, Germany \\ \faPhone: +49 3641 9-49786  | 
        \faAt: \href{mailto:johannes.zimmermann@uni-jena.de }{johannes.zimmermann@uni-jena.de}}
        
    \\






    % \cvitemshort{\href{http://lab42open.hcmr.gr/people/evangelospafilis/}{Pafilis Evangelos}}
    %     {Institute of Marine Biology, Biotechnology \& Aquaculture, Hellenic Centre for Marine Research, (IMBBC - HCMR
    % % {\href{https://imbbc.hcmr.gr/}{IMBBC HCMR}}
    %     ),
    %     Gournes, Pediados, P.O. Box 2214, Heraklion Crete 71003, Greece \\ \faPhone: +30 2810 337740 | 
    %     \faAt: \href{mailto:pafilis@hcmr.gr}{pafilis@hcmr.gr}}


    % \cvitemshort{\href{https://imbbc.hcmr.gr/user/kotoulas/}{Georgios Kotoulas}}
    %     {Institute of Marine Biology, Biotechnology \& Aquaculture, Hellenic Centre for Marine Research, (IMBBC - HCMR) \\
    %     Gournes, Pediados, P.O. Box 2214, Heraklion Crete 71003, Greece \\ \faPhone: +30 2810 337740 | 
    %     \faAt: \href{mailto:kotoulas@hcmr.gr}{kotoulas@hcmr.gr}} 


    \cvitemshort{\href{https://imbbc.hcmr.gr/user/kotoulas/}{Γεώργιος Κοτούλας}}
        {Ινστιτούτο Θαλάσσιας Βιολογίας, Βιοτεχνολογίας \& Ιχθυοκαλλιεργιών, Ελληνικό Κέντρο Θαλάσσιας Βιολογίας, (ΙΘΑΒΒΥΚ - ΕΛΚΕΘΕ) \\
        Γούρνες, Πεδιάδος, Ηράκλειο, Κρήτη 71003, Ελλάδα \\ \faPhone: +30 2810 337740 | 
        \faAt: \href{mailto:kotoulas@hcmr.gr}{kotoulas@hcmr.gr}} 


    % \cvitemshort{\href{https://ladoukakis.weebly.com/}{Emmanuel D. Ladoukakis}}
    %     {Biology Department, University of Crete,
    %     Voutes University Campus, Iraklio, GR 70013 Biology building, 3rd floor, Office: 316a 
    %     \\ 
    %     \faPhone: +30 2810394067 | 
    %     \faAt: \href{mailto:ladoukakis@biology.uoc.gr}{ladoukakis@biology.uoc.gr}}
      
    % \\ 

    
    % \cvitemshort{\href{https://who.paris.inria.fr/Elias.Tsigaridas/}{Elias Tsigaridas}}
    %     {Institut de Mathématiques de Jussieu - Paris Rive Gauche, Sorbonne Université - Campus Pierre et Marie Curie, Case courrier 247, 4, place Jussieu, 75252, Paris Cedex 05, France \\ 
    %     \faPhone: +33 014427 8539 | \faAt: \href{mailto: elias.tsigaridas@inria.fr}{elias.tsigaridas@inria.fr}}

    % \\



    % \cvitemshort{\href{http://computational-genomics.weebly.com/}{Christoforos Nikolaou}}
    %     {Department of Biology, University of Crete, Gennimata 34,71305, Heraklion, Crete, Greece \\ 
    %     \faAt: \href{mailto:cnikolaou@fleming.gr}{cnikolaou@fleming.gr}}


   \end{cvtable}

\makebacksidebar


\end{document}
